\chapter{2014年四川卷（文）}

\section{单选题}

\begin{problem}
已知集合 \(A = \{x \mid (x + 1)(x - 2) \leqslant 0\}\)，集合 \(B\) 为整数集，则 \(A \cap B =\)（\quad）

\choices
    {\( \{-1,0\} \)}
    {\(\{0, 1\}\)}
    {\(\{-2, -1, 0, 1\}\)}
    {\(\{-1, 0, 1, 2\}\)}
\begin{answer}
D
\end{answer}

\begin{solution}
由 \((x+1)(x-2)\leqslant0\) 的两个根为 \(-1\)，\(2\)，且二次项系数为正，得 \(A=[-1,2]\). 又 \(B=\Z\)，所以 \(A\cap B=\{-1,0,1,2\}\)，故选 D.
\end{solution}
\end{problem}

\begin{problem}
在“世界读书日”前夕，为了了解某地 5000 名居民某天的阅读时间，从中抽取了 200 名居民的阅读时间进行统计分析. 在这个问题中，5000 名居民的阅读时间的全体是（\quad）

\choices
    {总体}
    {个体}
    {样本的容量}
    {从总体中抽取的一个样本}
\begin{answer}
A
\end{answer}

\begin{solution}
统计对象是这 \(5000\) 名居民某天的阅读时间，研究中的全部个体组成总体；其中任意一名居民的阅读时间是一个个体，抽取的 \(200\) 个阅读时间构成样本，\(200\) 是样本容量. 因此题目所说的全体是总体，故选 A.
\end{solution}
\end{problem}

\begin{problem}
为了得到函数 \(y = \sin (x + 1)\) 的图象，只需把函数 \(y = \sin x\) 的图象上所有的点（\quad）

\choices
    {向左平行移动 1 个单位长度}
    {向右平行移动 1 个单位长度}
    {向左平行移动 \(\pi\) 个单位长度}
    {向右平行移动 \(\pi\) 个单位长度}
\begin{answer}
A
\end{answer}

\begin{solution}
把 \(y=\sin x\) 的图象向左平移 \(1\) 个单位后，图象上原来横坐标为 \(x+1\) 的点移动到横坐标为 \(x\) 的位置，函数值变为 \(\sin(x+1)\). 因此所得函数为 \(y=\sin(x+1)\)，故选 A.
\end{solution}
\end{problem}

\begin{problem}
某三棱锥的侧视图，俯视图如图所示，则该三棱锥的体积是（\quad）

\bitmapfigure[max width=0.44\linewidth]{img/2014/sichuan_liberal/q4_fig1.png}

\choices
    {3}
    {2}
    {\(\sqrt{3}\)}
    {1}
\begin{answer}
D
\end{answer}

\begin{solution}
由侧视图可知三棱锥的高为等腰三角形的高. 该等腰三角形的两腰长为 \(2\)，底边为 \(1+1=2\)，故
\[
h=\sqrt{2^{2}-1^{2}}=\sqrt{3}
\]
由俯视图可知底面为边长为 \(2\) 的等边三角形，底面积为 \(S=\frac12\times2\times\sqrt3=\sqrt3\). 所以三棱锥的体积为 \(V=\frac13Sh=\frac13\times\sqrt3\times\sqrt3=1\)，故选 D.
\end{solution}
\end{problem}

\begin{problem}
若 \(a > b > 0\)，\(c < d < 0\)，则一定有（\quad）

\choices
    {\(\frac{a}{c} > \frac{b}{d}\)}
    {\(\frac{a}{c} < \frac{b}{d}\)}
    {\(\frac{a}{d} > \frac{b}{c}\)}
    {\(\frac{a}{d} < \frac{b}{c}\)}
\begin{answer}
D
\end{answer}

\begin{solution}
令 \(c=-u\)，\(d=-v\)，则 \(u>v>0\). 比较后两个选项，因 \(cd>0\)，有
\[
\frac{a}{d}<\frac{b}{c}\iff ac<bd\iff au>bv
\]
而 \(a>b>0\)，\(u>v>0\)，所以 \(au>bv\)，从而 \(\frac{a}{d}<\frac{b}{c}\) 恒成立，故选 D.
\end{solution}
\end{problem}

\begin{problem}
执行如图所示的程序框图，如果输入的 \(x\)，\(y \in \R\)，则输出的 \(S\) 的最大值为（\quad）

\texfigure[float=inline,max width=5.8cm]{img_repaint/2014/sichuan_liberal/q6_fig1.tex}

\choices
    {0}
    {1}
    {2}
    {3}
\begin{answer}
C
\end{answer}

\begin{solution}
当输入的 \(x\)，\(y\) 满足 \(x\geqslant0\)，\(y\geqslant0\)，\(x+y\leqslant1\) 时，程序输出 \(S=2x+y\). 该可行域的顶点为 \((0,0)\)，\((1,0)\)，\((0,1)\). 在三个顶点处，\(2x+y\) 的值分别为 \(0\)，\(2\)，\(1\)，故线性函数在可行域上的最大值为 \(2\). 不满足条件时程序输出 \(S=1\)，不能超过 \(2\). 因此输出的 \(S\) 的最大值为 \(2\)，故选 C.
\end{solution}
\end{problem}

\begin{problem}
已知 \(b > 0\)，\(\log_5 b = a\)，\(\lg b = c\)，\(5^{d} = 10\)，则下列等式一定成立的是（\quad）

\choices
    {\(d = ac\)}
    {\(a = cd\)}
    {\(c = ad\)}
    {\(d = a + c\)}
\begin{answer}
B
\end{answer}

\begin{solution}
由 \(5^{d}=10\) 得 \(d=\log_{5}10\). 换底可得
\[
c=\lg b=\frac{\log_{5}b}{\log_{5}10}=\frac{a}{d}
\]
于是 \(a=cd\)，故选 B.
\end{solution}
\end{problem}

\begin{problem}
如图，从气球 \(A\) 上测得正前方的河流的两岸 \(B\)，\(C\) 的俯角分别为 \(75^{\circ}\)，\(30^{\circ}\)，此时气球的高是 \(60 \mathrm{~m}\)，则河流的宽度 \(BC\) 等于（\quad）

\bitmapfigure[float=inline,max width=0.42\linewidth]{img/2014/sichuan_liberal/q8_fig1.png}

\choices
    {\(240(\sqrt{3} - 1)\mathrm{\,m}\)}
    {\(180(\sqrt{2} - 1)\mathrm{\,m}\)}
    {\(120(\sqrt{3} - 1)\mathrm{\,m}\)}
    {\(30(\sqrt{3} + 1)\mathrm{\,m}\)}
\begin{answer}
C
\end{answer}

\begin{solution}
设 \(A\) 在地面上的垂足为 \(D\)，则 \(AD=60\mathrm{~m}\). 由俯角等于相应的仰角，在直角三角形 \(ADB\)，\(ADC\) 中分别有
\(DB=60\tan15^{\circ}\)，\(DC=60\tan60^{\circ}\).
又 \(\tan15^{\circ}=\tan(45^{\circ}-30^{\circ})=2-\sqrt3\)，故
\[
BC=DC-DB=60\sqrt3-60(2-\sqrt3)=120(\sqrt3-1)\mathrm{~m}
\]
因此河流宽度对应选项 C.
\end{solution}
\end{problem}

\begin{problem}
设 \(m \in \R\)，过定点 \(A\) 的动直线 \(x + my = 0\) 和过定点 \(B\) 的动直线 \(mx - y - m + 3 = 0\) 交于点 \(P(x, y)\). 则 \(|PA| + |PB|\) 的取值范围是（\quad）

\choices
    {\([\sqrt{5}, 2\sqrt{5}]\)}
    {\([\sqrt{10}, 2\sqrt{5}]\)}
    {\([\sqrt{10}, 4\sqrt{5}]\)}
    {\([2\sqrt{5}, 4\sqrt{5}]\)}
\begin{answer}
B
\end{answer}

\begin{solution}
两条动直线分别恒过 \(A=(0,0)\) 和 \(B=(1,3)\). 它们的法向量分别为 \((1,m)\) 和 \((m,-1)\)，数量积为 \(m-m=0\)，所以两直线始终垂直，即 \(\angle APB=90^{\circ}\). 因而 \(P\) 在以 \(AB\) 为直径的圆上，且
\(|AB|=\sqrt{(1-0)^2+(3-0)^2}=\sqrt{10}\)，\(|PA|^2+|PB|^2=10\).
由基本不等式，
\[
\sqrt{10}\leqslant |PA|+|PB|\leqslant\sqrt{2\left(|PA|^2+|PB|^2\right)}=2\sqrt5
\]
联立两条直线方程可得
\[
P=\left(\frac{m(m-3)}{m^2+1},\frac{3-m}{m^2+1}\right)
\]
当 \(m=3\) 或 \(m=-\frac13\) 时，分别有 \(P=A\) 或 \(P=B\)，下界可以取得；当 \(m=\frac12\) 或 \(m=-2\) 时，分别有 \(P=(-1,2)\) 或 \(P=(2,1)\)，均满足 \(|PA|=|PB|=\sqrt5\)，上界也可以取得. 故取值范围为 \([\sqrt{10},2\sqrt5]\)，选 B.
\end{solution}
\end{problem}

\begin{problem}
已知 \(F\) 是抛物线 \(y^{2} = x\) 的焦点，点 \(A\)，\(B\) 在该抛物线上且位于 \(x\) 轴的两侧， \(\overrightarrow{OA} \cdot \overrightarrow{OB} = 2\)（其中 \(O\) 为坐标原点），则 \(\triangle ABO\) 与 \(\triangle AFO\) 面积之和的最小值是（\quad）

\choices
    {2}
    {3}
    {\(\frac{17\sqrt{2}}{8}\)}
    {\(\sqrt{10}\)}
\begin{answer}
B
\end{answer}

\begin{solution}
设 \(A=(u^2,u)\)，\(B=(v^2,v)\)，其中 \(u>0>v\). 由 \(\overrightarrow{OA}\cdot\overrightarrow{OB}=2\)，得
\[
u^2v^2+uv=2
\]
令 \(t=uv<0\)，则 \(t^2+t-2=0\)，所以 \(t=-2\)，即 \(uv=-2\). 抛物线 \(y^2=x\) 的焦点为 \(F\left(\frac14,0\right)\). 因此
\(S_{\triangle ABO}=\frac12|u^2v-v^2u|=u-v\)，\(S_{\triangle AFO}=\frac12\cdot\frac14\cdot u=\frac{u}{8}\).
又由 \(v=-\frac2u\)，有
\[
S_{\triangle ABO}+S_{\triangle AFO}=\frac98u+\frac2u\geqslant2\sqrt{\frac98u\cdot\frac2u}=3
\]
等号条件为 \(u=\frac43\)，此时 \(v=-\frac32\)，满足 \(u>0>v\)，所以等号确实可以取得. 故最小值为 \(3\)，选 B.
\end{solution}
\end{problem}

\section{填空题}

\begin{problem}
双曲线 \(\frac{x^{2}}{4} - y^{2} = 1\) 的离心率等于\fillinblank{}.
\begin{answer}
\(\frac{\sqrt{5}}{2}\)
\end{answer}

\begin{solution}
双曲线的标准方程为 \(\frac{x^2}{a^2}-\frac{y^2}{b^2}=1\)，其中 \(a^2=4\)，\(b^2=1\). 由 \(c^2=a^2+b^2=5\) 得 \(c=\sqrt5\)，所以离心率为 \(e=\frac{c}{a}=\frac{\sqrt5}{2}\).
\end{solution}
\end{problem}

\begin{problem}
复数 \(\frac{2 - 2\i}{1 + \i} =\)\fillinblank{}.
\begin{answer}
\(-2\i\)
\end{answer}

\begin{solution}
分子、分母同乘以 \(1-\i\)，得
\[
\frac{2-2\i}{1+\i}=\frac{(2-2\i)(1-\i)}{(1+\i)(1-\i)}=\frac{-4\i}{2}=-2\i
\]
所以填 \(-2\i\).
\end{solution}
\end{problem}

\begin{problem}
设 \(f(x)\) 是定义在 \(\R\) 上的周期为 2 的函数，当 \(x \in [-1,1)\) 时，\(f(x) = \begin{cases}
-4x^{2} + 2, & -1 \leqslant x < 0,\\
    x, & 0 \leqslant x < 1,
\end{cases}\) 则 \(f\left(\frac{3}{2}\right) =\)\fillinblank{}.
\begin{answer}
1
\end{answer}

\begin{solution}
函数周期为 \(2\)，所以 \(f\left(\frac32\right)=f\left(\frac32-2\right)=f\left(-\frac12\right)\). 由于 \(-\frac12\in[-1,0)\)，代入分段表达式得
\[
f\left(-\frac12\right)=-4\left(-\frac12\right)^2+2=1
\]
故填 \(1\).
\end{solution}
\end{problem}

\begin{problem}
平面向量 \(\bs{a} = (1, 2)\)，\(\bs{b} = (4, 2)\)，\(\bs{c} = m\bs{a} + \bs{b}\)（\(m \in \R\)），且 \(\bs{c}\) 与 \(\bs{a}\) 的夹角等于 \(\bs{c}\) 与 \(\bs{b}\) 的夹角，则 \(m =\)\fillinblank{}.
\begin{answer}
2
\end{answer}

\begin{solution}
由 \(\bs{c}=m\bs{a}+\bs{b}\)，得 \(\bs{c}=(m+4,2m+2)\). 因为 \(|\bs{a}|=\sqrt5\)，\(|\bs{b}|=2\sqrt5\)，且两夹角相等，所以
\[
\frac{\bs{c}\cdot\bs{a}}{|\bs{c}|\,|\bs{a}|}=\frac{\bs{c}\cdot\bs{b}}{|\bs{c}|\,|\bs{b}|}
\]
即 \(2(\bs{c}\cdot\bs{a})=\bs{c}\cdot\bs{b}\). 计算得
\(\bs{c}\cdot\bs{a}=5m+8\)，\(qquad \bs{c}\cdot\bs{b}=8m+20\).
故 \(2(5m+8)=8m+20\)，解得 \(m=2\).
\end{solution}
\end{problem}

\begin{problem}
以 \(A\) 表示值域为 \(\R\) 的函数组成的集合，\(B\) 表示具有如下性质的函数 \(\varphi(x)\) 组成的集合：对于函数 \(\varphi(x)\)，存在一个正数 \(M\)，使得函数 \(\varphi(x)\) 的值域包含于区间 \([-M, M]\). 例如，当 \(\varphi_{1}(x) = x^{3}\)，\(\varphi_{2}(x) = \sin x\) 时，\(\varphi_{1}(x) \in A\)，\(\varphi_{2}(x) \in B\). 现有如下命题：

\begin{circlelist}
\item 设函数 \(f(x)\) 的定义域为 \(D\)，则“\(f(x) \in A\)”的充要条件是“\(\forall b \in \R\)，\(\exists a \in D\)，\(f(a)=b\)”；
\item 函数 \(f(x) \in B\) 的充要条件是 \(f(x)\) 有最大值和最小值；
\item 若函数 \(f(x)\)，\(g(x)\) 的定义域相同，且 \(f(x) \in A\)，\(g(x) \in B\)，则 \(f(x)+g(x) \notin B\)；
\item 若函数 \(f(x) = a\ln(x + 2) + \frac{x}{x^{2} + 1}\)（\(x > -2\)，\(a \in \R\)）有最大值，则 \(f(x) \in B\).
\end{circlelist}

其中的真命题有\fillinblank{}.（写出所有真命题的序号）
\begin{answer}
\(\circled{1}\circled{3}\circled{4}\)
\end{answer}

\begin{solution}
命题 1 直接是值域为 \(\R\) 的定义，因此命题 1 为真.

命题 2 中，函数有最大值和最小值可以推出值域有界，但有界函数不一定取得最大值和最小值. 例如 \(f(x)=\arctan x\) 的值域为 \(\left(-\frac\pi2,\frac\pi2\right)\)，有界却没有最大值和最小值，因此命题 2 为假.

命题 3 中，设 \(g(x)\) 的值域包含于 \([-M,M]\). 对任意正数 \(K\)，由于 \(f(x)\) 的值域为 \(\R\)，存在 \(x\) 使 \(f(x)>K+M\)，于是 \(f(x)+g(x)>K\). 因而 \(f+g\) 无上界，不能属于 \(B\)，命题 3 为真.

命题 4 中，若 \(a>0\)，则 \(a\ln(x+2)\to+\infty\)（\(x\to+\infty\)），函数不可能有最大值；若 \(a<0\)，则 \(a\ln(x+2)\to+\infty\)（\(x\to-2^+\)），函数同样不可能有最大值. 因此函数有最大值时只能有 \(a=0\)，此时 \(f(x)=\frac{x}{x^2+1}\)，且 \(\left|\frac{x}{x^2+1}\right|\leqslant\frac12\)，所以 \(f(x)\in B\). 命题 4 为真.

所以真命题的序号为 \(1\)，\(3\)，\(4\).
\end{solution}
\end{problem}

\section{解答题}

\begin{problem}
一个盒子里装有三张卡片，分别标记有数字 1， 2， 3，这三张卡片除标记的数字外完全相同. 随机有放回地抽取 3 次，每次抽取 1 张，将抽取的卡片上的数字依次记为 \(a\)，\(b\)，\(c\).

\begin{enumerate}
\item 求“抽取的卡片上的数字满足 \(a + b = c\)”的概率；
\item 求“抽取的卡片上的数字 \(a\)，\(b\)，\(c\) 不完全相同”的概率.
\end{enumerate}

\begin{solution}
\begin{enumerate}
\item 每次抽取都有 \(3\) 种等可能结果，三次抽取的有序结果共有 \(3^3=27\) 个. 满足 \(a+b=c\) 时，若 \(a=b=1\)，则 \(c=2\)，有 \(1\) 个结果；若 \((a,b)=(1,2)\) 或 \((2,1)\)，则 \(c=3\)，有 \(2\) 个结果；其他情形均不满足条件. 因此所求概率为 \(\frac{1+2}{27}=\frac19\).
\item 三次抽取结果等可能. \(a\)，\(b\)，\(c\) 完全相同的结果为 \((1,1,1)\)，\((2,2,2)\)，\((3,3,3)\)，共有 \(3\) 个，所以三者完全相同的概率为 \(\frac3{27}=\frac19\). 因此三者不完全相同的概率为 \(1-\frac19=\frac89\).
\end{enumerate}
\end{solution}
\end{problem}

\begin{problem}
已知函数 \(f(x) = \sin \left(3x + \frac{\pi}{4}\right)\).

\begin{enumerate}
\item 求 \(f(x)\) 的单调递增区间；
\item 若 \(\alpha\) 是第二象限角，\(f\left(\frac{\alpha}{3}\right) = \frac{4}{5}\cos \left(\alpha + \frac{\pi}{4}\right)\cos 2\alpha\)，求 \(\cos \alpha - \sin \alpha\) 的值.
\end{enumerate}

\begin{solution}
\begin{enumerate}
\item \(y=\sin t\) 在区间 \(\left[-\frac\pi2+2k\pi,\frac\pi2+2k\pi\right]\) 上单调递增. 令 \(t=3x+\frac\pi4\)，则
\(-\frac\pi2+2k\pi\leqslant3x+\frac\pi4\leqslant\frac\pi2+2k\pi\)，\(k\in\Z\).
解得 \(\left[\frac{2k\pi}{3}-\frac\pi4,\frac{2k\pi}{3}+\frac\pi{12}\right]\)，\(k\in\Z\)，这就是 \(f(x)\) 的单调递增区间.
\item 因为 \(f\left(\frac\alpha3\right)=\sin\left(\alpha+\frac\pi4\right)\)，设 \(u=\sin\alpha+\cos\alpha\)，\(v=\cos\alpha-\sin\alpha\). 已知等式化为
\[
\frac{u}{\sqrt2}=\frac45\frac{v}{\sqrt2}\cos2\alpha
\]
又 \(\cos2\alpha=(\cos\alpha-\sin\alpha)(\cos\alpha+\sin\alpha)=uv\)，所以 \(u=\frac45uv^2\). 若 \(u=0\)，则 \(v=-\sqrt2\). 若 \(u\neq0\)，则 \(v^2=\frac54\). 由于 \(\alpha\) 是第二象限角，\(\cos\alpha-\sin\alpha<0\)，故 \(v=-\frac{\sqrt5}{2}\). 因此 \(\cos\alpha-\sin\alpha=-\sqrt2\) 或 \(-\frac{\sqrt5}{2}\).
\end{enumerate}
\end{solution}
\end{problem}

\begin{problem}
在如图所示的多面体中，四边形 \(ABB_{1}A_{1}\) 和 \(ACC_{1}A_{1}\) 都为矩形.

\begin{enumerate}
\item 若 \(AC \perp BC\)，证明：直线 \(BC \perp\) 平面 \(ACC_{1}A_{1}\)；
\item 设 \(D\)，\(E\) 分别是线段 \(BC\)，\(CC_{1}\) 的中点，在线段 \(AB\) 上是否存在一点 \(M\)，使直线 \(DE \parallel\) 平面 \(A_{1}MC\)？请证明你的结论.
\end{enumerate}

\bitmapfigure[max width=0.46\linewidth]{img/2014/sichuan_liberal/q18_fig1.png}

\begin{solution}
\begin{enumerate}
\item 设 \(\bs{u}=\overrightarrow{AA_{1}}\)，\(\bs{v}=\overrightarrow{AB}\)，\(\bs{w}=\overrightarrow{AC}\). 由两个四边形为矩形，得 \(\bs{u}\cdot\bs{v}=0\)，\(\bs{u}\cdot\bs{w}=0\). 又 \(\overrightarrow{BC}=\bs{w}-\bs{v}\)，所以 \(\overrightarrow{BC}\cdot\bs{u}=0\)，即 \(BC\perp AA_{1}\). 已知 \(AC\perp BC\)，而 \(AC\) 与 \(AA_{1}\) 是平面 \(ACC_{1}A_{1}\) 内相交的两条直线，因此 \(BC\perp\) 平面 \(ACC_{1}A_{1}\).
\item 取 \(M\) 为 \(AB\) 的中点. 用上面的向量记号，有
\(\overrightarrow{DE}=\frac12\left(\bs{u}+\bs{w}-\bs{v}\right)\)，\(\overrightarrow{A_{1}M}=\frac12\bs{v}-\bs{u}\)，\(\overrightarrow{A_{1}C}=\bs{w}-\bs{u}\).
从而
\[
\overrightarrow{DE}=-\overrightarrow{A_{1}M}+\frac12\overrightarrow{A_{1}C}
\]
右端是平面 \(A_{1}MC\) 内两条相交直线方向向量的线性组合，所以直线 \(DE\) 的方向平行于平面 \(A_{1}MC\). 又因 \(ABC\) 是三角形，且 \(\bs{u}\) 同时垂直于 \(\bs{v}\)，\(\bs{w}\)，故 \(\bs{u}\)，\(\bs{v}\)，\(\bs{w}\) 线性无关. 若 \(D\) 在平面 \(A_{1}MC\) 内，则存在实数 \(\lambda\)，\(\mu\)，使
\[
\overrightarrow{A_{1}D}=\frac12(\bs{v}+\bs{w})-\bs{u}=\lambda\left(\frac12\bs{v}-\bs{u}\right)+\mu(\bs{w}-\bs{u})
\]
比较 \(\bs{v}\)，\(\bs{w}\)，\(\bs{u}\) 的系数，分别得到 \(\lambda=1\)，\(\mu=\frac12\)，\(\lambda+\mu=1\)，矛盾. 因此 \(D\) 不在平面 \(A_{1}MC\) 内，故 \(DE\parallel\) 平面 \(A_{1}MC\). 因此存在这样的点 \(M\)，且它是线段 \(AB\) 的中点.
\end{enumerate}
\end{solution}
\end{problem}

\begin{problem}
设等差数列 \(\{a_{n}\}\) 的公差为 \(d\)，点 \((a_{n}, b_{n})\) 在函数 \(f(x) = 2^{x}\) 的图象上（\(n \in \N^{*}\)).

\begin{enumerate}
\item 证明：数列 \(\{b_{n}\}\) 为等比数列；
\item 若 \(a_{1} = 1\)，函数 \(f(x)\) 的图象在点 \((a_{2}, b_{2})\) 处的切线在 \(x\) 轴上的截距为 \(2 - \frac{1}{\ln 2}\)，求数列 \(\{a_{n} b_{n}^{2}\}\) 的前 \(n\) 项和 \(S_{n}\).
\end{enumerate}

\begin{solution}
\begin{enumerate}
\item 因为点 \((a_n,b_n)\) 在 \(y=2^x\) 的图象上，所以 \(b_n=2^{a_n}\). 又 \(a_{n+1}=a_n+d\)，故
\[
\frac{b_{n+1}}{b_n}=\frac{2^{a_n+d}}{2^{a_n}}=2^d
\]
这个比值与 \(n\) 无关，所以 \(\{b_n\}\) 是以 \(2^d\) 为公比的等比数列.
\item 函数 \(f(x)=2^x\) 的导数为 \(f'(x)=2^x\ln2\). 在点 \((a_2,b_2)\) 处的切线方程为
\[
y-b_2=b_2\ln2\left(x-a_2\right)
\]
令 \(y=0\)，得该切线在 \(x\) 轴上的截距为 \(a_2-\frac1{\ln2}\). 由题设及 \(a_1=1\)，得 \(a_2=2\)，从而 \(d=1\)，\(a_n=n\)，\(b_n=2^n\). 因此
\[
S_n=\sum_{k=1}^n a_kb_k^2=\sum_{k=1}^n k4^k
\]
设上式为 \(S_n\)，则 \(4S_n=\sum_{k=2}^{n+1}(k-1)4^k\)，相减得
\[
3S_n=n4^{n+1}-\sum_{k=1}^n4^k=n4^{n+1}-\frac{4^{n+1}-4}{3}
\]
所以 \(S_n=\frac{(3n-1)4^{n+1}+4}{9}\).
\end{enumerate}
\end{solution}
\end{problem}

\begin{problem}
已知椭圆 \(C: \frac{x^2}{a^2} + \frac{y^2}{b^2} = 1\)（\(a > b > 0\)）的左焦点为 \(F(-2,0)\)，离心率为 \(\frac{\sqrt{6}}{3}\).

\begin{enumerate}
\item 求椭圆 \(C\) 的标准方程；
\item 设 \(O\) 为坐标原点，\(T\) 为直线 \(x = -3\) 上一点，过 \(F\) 作 \(TF\) 的垂线交椭圆于 \(P\)，\(Q\). 当四边形 \(OPTQ\) 是平行四边形时，求四边形 \(OPTQ\) 的面积.
\end{enumerate}

\begin{solution}
\begin{enumerate}
\item 设椭圆焦距参数为 \(c\). 由左焦点 \(F(-2,0)\) 得 \(c=2\)，再由离心率 \(e=\frac ca=\frac{\sqrt6}{3}\)，得 \(a=\sqrt6\). 因为 \(b^2=a^2-c^2=6-4=2\)，所以椭圆的标准方程为 \(\frac{x^2}{6}+\frac{y^2}{2}=1\).
\item 设 \(T=(-3,m)\). 向量 \(\overrightarrow{TF}=(1,-m)\)，过 \(F\) 且垂直于 \(TF\) 的直线可写为 \(x=my-2\). 设 \(P=(x_1,y_1)\)，\(Q=(x_2,y_2)\)，代入椭圆方程得
\[
(m^2+3)y^2-4my-2=0
\]
由韦达定理，
\(y_1+y_2=\frac{4m}{m^2+3}\)，\(y_1y_2=-\frac2{m^2+3}\).
又 \(x_i=my_i-2\)，所以 \(x_1+x_2=\frac{4m^2}{m^2+3}-4=-\frac{12}{m^2+3}\). 四边形 \(OPTQ\) 为平行四边形时，对角线互相平分，即 \(P+Q=O+T=T\)，故
\(\frac{4m}{m^2+3}=m\)，\(-\frac{12}{m^2+3}=-3\).
第二个等式给出 \(m^2=1\)，这两个值也满足第一个等式，因此 \(m=\pm1\).

取 \(m=1\)，则 \(T=(-3,1)\)，直线 \(PQ\) 为 \(y=x+2\)，与椭圆的两个交点为
\(P\left(\frac{-3-\sqrt3}{2},\frac{1-\sqrt3}{2}\right)\)，\(Q\left(\frac{-3+\sqrt3}{2},\frac{1+\sqrt3}{2}\right)\).
平行四边形的两条对角线为 \(OT\) 与 \(PQ\)，故面积为
\[
\frac12\left|\det\left((-3,1),(\sqrt3,\sqrt3)\right)\right|=2\sqrt3
\]
当 \(m=-1\) 时图形关于 \(x\) 轴对称，面积相同. 因此四边形 \(OPTQ\) 的面积为 \(2\sqrt3\).
\end{enumerate}
\end{solution}
\end{problem}

\begin{problem}
已知函数 \(f(x) = \e^{x} - ax^{2} - bx - 1\)，其中 \(a\)，\(b \in \R\)，\(\e = 2.71828\cdots\) 为自然对数的底数.

\begin{enumerate}
\item 设 \(g(x)\) 是函数 \(f(x)\) 的导函数，求函数 \(g(x)\) 在区间 \([0,1]\) 上的最小值；
\item 若 \(f(1)=0\)，函数 \(f(x)\) 在区间 \((0,1)\) 内有零点，证明：\(\e - 2 < a < 1\).
\end{enumerate}

\begin{solution}
\begin{enumerate}
\item \(g(x)=f'(x)=\e^x-2ax-b\)，且 \(g'(x)=\e^x-2a\) 严格递增. 若 \(a\leqslant\frac12\)，则 \(g'(x)\geqslant0\)（\(0\leqslant x\leqslant1\)），所以最小值为 \(g(0)=1-b\). 若 \(\frac12<a<\frac\e2\)，则 \(g'(x)=0\) 的解 \(x_0=\ln(2a)\) 在 \((0,1)\) 内，故最小值为
\[
g(x_0)=2a-2a\ln(2a)-b
\]
若 \(a\geqslant\frac\e2\)，则 \(g'(x)\leqslant0\)（\(0\leqslant x\leqslant1\)），所以最小值为 \(g(1)=\e-2a-b\). 综上，最小值为
\[
\begin{cases}
1-b, & a\leqslant\frac12,\\
2a-2a\ln(2a)-b, & \frac12<a<\frac\e2,\\
\e-2a-b, & a\geqslant\frac\e2.
\end{cases}
\]
\item 由 \(f(1)=0\)，得 \(a+b=\e-1\)，即 \(b=\e-1-a\). 又 \(f(0)=0\)，设题设中 \((0,1)\) 内的零点为 \(x_0\). 对区间 \([0,x_0]\) 和 \([x_0,1]\) 分别使用罗尔定理，得 \(g=f'\) 在 \((0,1)\) 内有两个不同的零点.

由于 \(g'(x)=\e^x-2a\) 严格递增，且 \(g\) 在 \((0,1)\) 内有两个不同的零点，罗尔定理表明 \(g'\) 在这两个零点之间有一个零点. 因而 \(g\) 在该点左侧递减、右侧递增，两个零点之外函数值严格为正，从而 \(g(0)>0\)，\(g(1)>0\). 因此
\(g(0)=1-b=a+2-\e>0\)，\(g(1)=\e-2a-b=1-a>0\).
于是 \(\e-2<a<1\).
\end{enumerate}
\end{solution}
\end{problem}
